\documentclass[10pt,a4paper]{article}
\usepackage[margin=1.45cm]{geometry}
\usepackage{fontspec}
\usepackage{xeCJK}
\usepackage{array}
\usepackage{amsmath,amssymb,amsfonts}
\usepackage{booktabs}
\usepackage{enumitem}
\usepackage{hyperref}
\usepackage[table]{xcolor}
\usepackage{titlesec}

\hypersetup{colorlinks=true,linkcolor=black,urlcolor=blue,citecolor=black}
\setmainfont{Times New Roman}
\setCJKmainfont{Songti SC}
\setlength{\parindent}{0pt}
\setlength{\parskip}{2pt}
\setlist[itemize]{leftmargin=1.1em,itemsep=0.8pt,topsep=1.5pt}
\setlist[enumerate]{leftmargin=1.25em,itemsep=0.8pt,topsep=1.5pt}
\renewcommand{\arraystretch}{1.08}
\titleformat{\section}{\normalsize\bfseries}{\thesection}{0.55em}{}
\titlespacing*{\section}{0pt}{4pt}{1.5pt}
\titleformat{\subsection}{\normalsize\bfseries}{\thesubsection}{0.45em}{}
\titlespacing*{\subsection}{0pt}{3pt}{1pt}

\newcommand{\blank}[1][7em]{\colorbox{yellow!35}{\rule{0pt}{1.05em}\hspace{#1}}}
\newcommand{\eg}{E_{\mathrm{g}}}
\newcommand{\ehull}{E_{\mathrm{hull}}}
\newcommand{\form}{x_{\mathrm{form}}}
\newcommand{\struc}{x_{\mathrm{struc}}}
\newcommand{\unc}{u}
\newcommand{\stab}{S_{\mathrm{stab}}}
\newcommand{\domainf}{d_{\mathrm{form}}}
\newcommand{\domains}{d_{\mathrm{struc}}}
\newcommand{\utility}{U}
\newcommand{\gapscore}{\rho_{\mathrm{gap}}}
\newcommand{\uvscore}{\rho_{\mathrm{UV/WBG}}}
\newcommand{\dielscore}{\rho_{\mathrm{dielectric/2D}}}
\newcommand{\novel}{\rho_{\mathrm{novel}}}

\begin{document}
\vspace*{-34pt}
\begin{center}
    \textbf{CITY UNIVERSITY OF HONG KONG}\\
    \vspace{6pt}
    \underline{\textbf{Research Plan Template}}
\end{center}

\section{Title of the Proposed Research}
\begin{itemize}[leftmargin=15pt]
    \item English Project Title: \textbf{AI-Powered Boron Nitride Material Exploration}
    \item Chinese Project Title: 人工智能驱动的氮化硼材料探索
\end{itemize}

\section{Principal Investigator, Collaborator(s), and Duration}
\begin{itemize}[leftmargin=15pt]
    \item Principal investigator: Prof. Chen MA
    \item Position and affiliated department: Assistant Professor, Department of Computer Science, City University of Hong Kong
    \item Collaborator(s) / stakeholder contact: Shenzhen Zhongneng Tianyi Technology Co., Ltd. (to be confirmed by the PI)
    \item Duration of the project, including proposed start date and scheduled completion date:\\[-1pt]\blank[18em]
\end{itemize}

\section{Abstract of the Project}
Boron nitride (BN) materials are attractive for ultraviolet/wide-band-gap and dielectric/2D-support applications, yet AI-guided exploration is limited by sparse labels, uneven coverage of BN variants, and uncertain extrapolation from broad materials datasets. This project aims to build an AI-powered BN material exploration framework that turns public materials data into uncertainty-aware research decisions. The framework will define a bounded BN-centered design space, screen formula-only candidates with composition models, generate prototype-first structure hypotheses for selected formulas, and rank candidates using target-window, stability, uncertainty, domain-distance, and novelty criteria. Expected outputs include a curated BN dataset, benchmarked models, ranked candidate table, validation-ready structure package, stakeholder review package, and manuscript-ready analysis.

\section{Key Objectives}
The central objective is to build a BN-specific, uncertainty-aware candidate-prioritization pipeline, rather than a generic list of materials-ML tasks:
\[
  \mathcal{P}_{\mathrm{BN}}:
  (\mathcal{D},\mathcal{C}_{\mathrm{BN}})\rightarrow
  \{(c,s,\hat{\mathbf{y}},\hat{\unc},a)\}_{k}.
\]

\begin{enumerate}
  \item \textbf{Define a bounded BN exploration space.}
  \[
    \mathcal{C}_{\mathrm{BN}}=
    \mathcal{C}_{\mathrm{core}}\cup\mathcal{C}_{\mathrm{ext}}\cup
    \mathcal{C}_{\mathrm{explore}},
  \]
  where $\mathcal{C}_{\mathrm{core}}$ covers BN polymorphs and doped BN,
  $\mathcal{C}_{\mathrm{ext}}$ covers BCN, BC$_2$N, and SiBN motifs, and
  $\mathcal{C}_{\mathrm{explore}}$ is restricted to charge-neutral B/N/C/Si-centered substitutions that map to small BN-family prototypes:
  \[
    \mathcal{C}_{\mathrm{explore}}=
    \{c:q(c)=0,\;N_{\mathrm{atom}}(c)\le 12,\;\pi(c)\in\Pi_{\mathrm{BN}}\}.
  \]
  \textit{Output:} a provenance-labelled dataset and candidate design space.

  \item \textbf{Develop two-stage property and uncertainty models.}
  Formula-only candidates will first use composition descriptors and composition-based models; only candidates passing this screen will enter structure hypothesis generation:
  \[
    f_{\theta}^{(1)}:\form\rightarrow(\hat{\eg},\hat{\stab},\hat{\unc}),
    \quad
    f_{\theta}^{(2)}:(\form,\struc)\rightarrow
    (\hat{\eg},\hat{\ehull},\hat{\boldsymbol{\epsilon}},\hat{\unc}).
  \]
  \textit{Output:} benchmarked formula-only and structure-resolved models.

  \item \textbf{Prioritize candidates under two application tracks.}
  The project focuses on two tracks only: UV/wide-band-gap candidates and dielectric/2D-support candidates. Candidate utility is
  \[
    \utility_t(c,s)=w_1\rho_t(c,s)+w_2\hat{\stab}(c,s)-w_3\hat{\unc}(c,s)
    -w_4 d_t(c,s)+w_5\novel(c),\quad
    t\in\{\mathrm{UV/WBG},\mathrm{dielectric/2D}\}.
  \]
  All terms are normalized to $[0,1]$; $d_t$ is $\domainf$ at formula stage and $\domains$ at structure stage. Weights will be set by sensitivity analysis and PI review, not claimed as globally optimal.
  \textit{Output:} a ranked candidate table with predicted properties, uncertainty, novelty, application label, and action label.

  \item \textbf{Produce validation-ready structure hypotheses.}
  Priority formulas will be converted into candidate structures by prototype reuse, database-derived structure transfer, and controlled substitution before any generative route:
  \[
    c\mapsto\{s_{c,1},\ldots,s_{c,k}\}\mapsto
    \{\mathrm{sanity\ check},\mathrm{relaxation},\mathrm{property\ recheck}\}.
  \]
  \textit{Output:} CIF/POSCAR files, provenance records, predicted properties, uncertainty scores, and validation notes.
\end{enumerate}

\section{Research Methodology}
\subsection{Background of Research}
BN is a focused but scientifically rich materials family. h-BN and c-BN anchor the design space, while B-C-N, BC$_2$N, Si-B-N, doped BN, and controlled small-cell analogues define a broader yet bounded exploration set. The research question is therefore not ``can one model predict all materials,'' but:
\[
  \textit{which BN-related formulas and structures should be investigated next, and with what uncertainty?}
\]
The project will build on preliminary formula-level candidate-ranking components, including grouped evaluation, BN-focused diagnostics, uncertainty/support penalties, and structure-follow-up artifacts, and will turn them into a general AI-powered BN material exploration framework.

\subsection{WP1: Data Construction and Provenance Tracking}
The project will integrate public computational resources, including 2DMatPedia, JARVIS, Matbench, matminer, the Materials Project, and C2DB~\cite{zhou2019_2dmatpedia,choudhary2020_jarvis,dunn2020_matbench,ward2018_matminer,jain2013_mp,haastrup2018_c2db}. Each record is represented as
\[
  r_i=(c_i,s_i,\mathbf{y}_i,p_i),\qquad
  p_i=(\mathrm{source},\mathrm{method},\mathrm{version},\mathrm{provenance}),
\]
where $c_i$ is composition, $s_i$ is optional structure, and $\mathbf{y}_i$ stores available target properties. The BN-centered subset is
\[
  \mathcal{D}_{\mathrm{BN}}=\{r_i:\{B,N\}\subseteq\mathrm{elements}(c_i)\ \mathrm{or}\ c_i\in\mathcal{C}_{\mathrm{BN}}\}.
\]
Deduplication, dimensionality labels, target-property alignment, and source-version tracking will be completed before modelling. Incompatible labels will be retained with source indicators rather than merged as identical ground truth.

\subsection{WP2: Formula-Only and Structure-Resolved Models}
The modelling layer will compare descriptor baselines, composition neural baselines, and structure-resolved graph models. Candidate screening begins formula-only:
\[
  \form\rightarrow
  [\mathrm{composition\ descriptors},\mathrm{element\ fractions}]
  \rightarrow f_{\theta}^{(1)}.
\]
After selected formulas obtain structures, structure-resolved models such as CGCNN, MEGNet, ALIGNN, and CHGNet can be applied~\cite{xie2018_cgcnn,chen2019_megnet,choudhary2021_alignn,deng2023_chgnet}. The band-gap objective will be a target-window score rather than a monotonic ``larger is always better'' term:
  \[
    \gapscore(c)=
    \exp\!\left[-\frac{(\hat{\eg}(c)-E^*_{\mathrm{g,app}})^2}{2\sigma_{\mathrm{app}}^2}\right]
    \cdot \eta_{\mathrm{direct}}(c).
  \]
At formula-only stage, directness is treated as unknown unless direct/indirect labels are explicitly available; $\eta_{\mathrm{direct}}$ is otherwise applied only after structure-resolved scoring or set as a conservative neutral/penalty term. The application scores $\uvscore$ and $\dielscore$ are proxy scores derived from computed properties, dimensionality metadata, and rule-based target criteria, not direct experimental labels. For the dielectric/2D-support track, ranking will use band-gap lower bounds, dimensionality/provenance labels, dielectric tensors where available, stability proxy, and uncertainty/domain-distance penalties rather than measured device performance. Predictive uncertainty will be estimated by model ensembles and calibrated against validation errors~\cite{gal2016_dropout,lakshminarayanan2017_ensembles}:
\[
  \hat{\unc}_{\mathrm{cal}}(c)=
  \alpha\,\mathrm{Std}_{m\in\mathcal{M}}[f_m(c)]
  +\beta\,\widehat{\mathrm{err}}_{\mathrm{val}}(c).
\]
If structure-resolved models do not outperform robust composition baselines under BN-family-aware splits, their outputs will be treated as exploratory evidence only.

\subsection{WP3: Candidate Ranking and Decision Protocol}
Candidate ranking will use prediction, reliability, and actionability. Formula-stage and structure-stage domain distances are separated:
\[
  \domainf(c)=\min_i\|\phi_{\mathrm{form}}(c)-\phi_{\mathrm{form}}(c_i)\|_{\Sigma_f},
  \quad
  \domains(c,s)=\min_i\|\phi_{\mathrm{struc}}(c,s)-\phi_{\mathrm{struc}}(c_i,s_i)\|_{\Sigma_s}.
\]
Novelty is estimated by distance from known BN-family compounds; application relevance is assigned by matching candidates to the UV/WBG or dielectric/2D-support track. The operational protocol is
\[
  \mathrm{candidate\ pool}\rightarrow\mathrm{property\ prediction}\rightarrow
  \mathrm{uncertainty/domain\ estimation}\rightarrow\mathrm{filtering}\rightarrow
  a(c),
\]
where
\[
  a(c)\in\{\mathrm{control},\mathrm{priority},\mathrm{explore},\mathrm{hold}\}.
\]
Priority candidates proceed to structure handoff and validation; exploratory candidates are retained as high-risk families; control candidates serve as reference checks; hold candidates are deferred until additional evidence is available. This design supports active learning, where high-uncertainty but promising families can be reviewed rather than discarded~\cite{lookman2019_active}.

\subsection{WP4: Structure Handoff and Validation}
Selected formulas will first be mapped to structures by prototype reuse, structure transfer, and controlled substitution. Generative models such as CDVAE will be considered only when prototype-based routes are insufficient~\cite{xie2022_cdvae}. Materials-side validation will include charge-neutrality checks, structural sanity checks, geometry relaxation where feasible, energy-above-hull re-estimation,
\[
  \ehull(c,s)=e(c,s)-
  \min_{\lambda_q\ge0}\sum_q\lambda_q e(q),
  \quad \sum_q\lambda_q=1,\;\sum_q\lambda_q\mathbf{n}_q=\mathbf{n}_c,
\]
and comparison with known BN-family references. Bulk-like entries will use formation-energy or energy-above-hull proxies; 2D entries will additionally use dimensionality, database provenance, and available 2D stability indicators. Model evaluation will use random, grouped-by-formula, and BN-family-aware splits, with MAE, RMSE, $R^2$, calibration error, Kendall $\tau$, Spearman $\rho$, and top-$k$ rank stability.

\section{Significance of the Proposed Research and Major Deliverables}
The significance of this project is a BN-specific AI decision pipeline: a curated data layer, family-aware evaluation, calibrated uncertainty, application-aware ranking, and prototype-first structure handoff. This is relevant to ultraviolet materials, insulating/dielectric layers, and 2D heterostructure support materials~\cite{watanabe2004_hbn_uv,cassabois2016_hbn_indirect,dean2010_hbn_graphene}. Thermal properties are secondary relevance rather than a third project target.

\textbf{Stakeholder feedback.}
Shenzhen Zhongneng Tianyi Technology Co., Ltd. may provide application-level feedback if arranged by the PI or donor office, with checkpoints after formula-stage ranking and structure-hypothesis packaging. The stakeholder-facing package will contain predicted properties, uncertainty estimates, domain-distance flags, novelty labels, validation notes, and recommended action labels, not unverified discovery claims.

\textbf{Major deliverables:}
\begin{enumerate}
  \item \textbf{D1:} BN-centered dataset from at least three public sources, with formula, structure, target-property, dimensionality, and provenance fields.
  \item \textbf{D2:} Benchmark report comparing descriptor, composition, and structure-resolved models under random, grouped-by-formula, and BN-family-aware splits.
  \item \textbf{D3:} Ranked table of at least 30 candidate families, including formula, predicted band gap, stability proxy, uncertainty, domain distance, novelty score, application label, and recommended action.
  \item \textbf{D4:} Structure hypothesis package for 5--10 priority candidates, including CIF/POSCAR files, provenance record, predicted properties, uncertainty score, and validation notes.
  \item \textbf{D5:} Technical report, figures, stakeholder review package, lightweight demo interface, and manuscript-ready analysis.
\end{enumerate}
Indicative milestones, to be adjusted after the PI confirms duration, are: data/provenance construction; model benchmarking; uncertainty calibration and ranking; structure package, review package, report, and manuscript-ready figures.

\small
\renewcommand{\refname}{Selected References}
\begin{thebibliography}{99}
\setlength{\itemsep}{-1pt}
\bibitem{zhou2019_2dmatpedia} J. Zhou et al., ``2DMatPedia: an open computational database of two-dimensional materials from top-down and bottom-up approaches,'' \textit{Scientific Data}, 6, 86, 2019.
\bibitem{choudhary2020_jarvis} K. Choudhary et al., ``JARVIS-Tools: open-access software for atomistic data-driven materials design,'' \textit{npj Computational Materials}, 6, 173, 2020.
\bibitem{dunn2020_matbench} A. Dunn et al., ``Benchmarking materials property prediction methods: the Matbench test set and Automatminer reference algorithm,'' \textit{npj Computational Materials}, 6, 138, 2020.
\bibitem{ward2018_matminer} L. Ward et al., ``Matminer: An open source toolkit for materials data mining,'' \textit{Computational Materials Science}, 152, 60--69, 2018.
\bibitem{jain2013_mp} A. Jain et al., ``The Materials Project: a materials genome approach to accelerating materials innovation,'' \textit{APL Materials}, 1, 011002, 2013.
\bibitem{haastrup2018_c2db} S. Haastrup et al., ``The Computational 2D Materials Database,'' \textit{2D Materials}, 5, 042002, 2018.
\bibitem{xie2018_cgcnn} T. Xie and J. C. Grossman, ``Crystal graph convolutional neural networks for materials property prediction,'' \textit{Physical Review Letters}, 120, 145301, 2018.
\bibitem{chen2019_megnet} C. Chen et al., ``Graph networks as a universal machine learning framework for molecules and crystals,'' \textit{Chemistry of Materials}, 31, 3564--3572, 2019.
\bibitem{choudhary2021_alignn} K. Choudhary and B. DeCost, ``Atomistic line graph neural network for improved materials property predictions,'' \textit{npj Computational Materials}, 7, 185, 2021.
\bibitem{deng2023_chgnet} B. Deng et al., ``CHGNet as a pretrained universal neural network potential,'' \textit{Nature Machine Intelligence}, 5, 1031--1041, 2023.
\bibitem{gal2016_dropout} Y. Gal and Z. Ghahramani, ``Dropout as a Bayesian approximation: representing model uncertainty in deep learning,'' \textit{ICML}, 2016.
\bibitem{lakshminarayanan2017_ensembles} B. Lakshminarayanan, A. Pritzel and C. Blundell, ``Simple and scalable predictive uncertainty estimation using deep ensembles,'' \textit{NeurIPS}, 2017.
\bibitem{lookman2019_active} T. Lookman, P. V. Balachandran, D. Xue and R. Yuan, ``Active learning in materials science with emphasis on adaptive sampling using uncertainties for targeted design,'' \textit{npj Computational Materials}, 5, 21, 2019.
\bibitem{xie2022_cdvae} T. Xie et al., ``Crystal diffusion variational autoencoder for periodic material generation,'' \textit{ICLR}, 2022.
\bibitem{watanabe2004_hbn_uv} K. Watanabe, T. Taniguchi and H. Kanda, ``Direct-bandgap properties and ultraviolet lasing of hexagonal boron nitride,'' \textit{Nature Materials}, 3, 404--409, 2004.
\bibitem{cassabois2016_hbn_indirect} G. Cassabois, P. Valvin and B. Gil, ``Hexagonal boron nitride is an indirect bandgap semiconductor,'' \textit{Nature Photonics}, 10, 262--266, 2016.
\bibitem{dean2010_hbn_graphene} C. R. Dean et al., ``Boron nitride substrates for high-quality graphene electronics,'' \textit{Nature Nanotechnology}, 5, 722--726, 2010.
\end{thebibliography}
\normalsize

\section{Ethics and Safety}
The initial project is computational and data-driven. It uses public databases, open-source software, models, and generated candidate structures. Curated records, scripts, model configurations, and candidate-ranking outputs will retain source identifiers and version metadata for reproducibility and license-aware reuse. No human subjects, living animals, wet-lab synthesis, biological samples, or radiation work will be conducted in the initial computational phase. If later experimental synthesis or regulated laboratory procedures are introduced, the responsible experimental party will obtain the required safety approvals before that work begins.

\clearpage
\section*{Project Details}
\textbf{Project Title:} AI-Powered Boron Nitride Material Exploration

\textbf{PI (Dept):} Prof. Chen MA (Department of Computer Science)

\textbf{Budget Breakdown}
\begin{center}
\small
\begin{tabular}{>{\raggedright\arraybackslash}p{0.18\linewidth}>{\raggedright\arraybackslash}p{0.57\linewidth}>{\raggedright\arraybackslash}p{0.15\linewidth}}
\toprule
\textbf{Budget Item} &
\textbf{Detail breakdown (please itemize)} &
\textbf{Amount (HKD)}\\
\midrule
Staffing & Student helper / research assistant support for data cleaning, benchmark scripts, candidate-table verification, figure generation, documentation, and demo maintenance. & \blank\\
Equipment & GPU cloud hours for graph-model training and ensemble inference; storage for downloaded structures, generated CIF/POSCAR files, benchmark logs, and model checkpoints. & \blank\\
General expenses & Cloud credits, repository backup, software services, stakeholder-facing demo hosting, and data-management costs. & \blank\\
Conference & Registration and presentation materials for AI-for-science, materials informatics, or applied computing venues. & \blank\\
Overseas Travel & Optional PI-approved stakeholder visit, technical meeting, or conference travel for candidate-review discussion. & \blank\\
\midrule
 & \textbf{Total Funding} & \blank\\
\bottomrule
\end{tabular}
\end{center}

\end{document}
